% !TeX program = xelatex
% Based on the supplied Chinese ctexart + spconf two-column template.
\documentclass[UTF8,a4paper,zihao=5,fontset=windows]{ctexart}
\usepackage{spconf}
\usepackage{amsmath,amssymb,graphicx,booktabs,multirow}
\usepackage{tikz}
\usetikzlibrary{arrows.meta,positioning,calc}
\usepackage{cite}
\usepackage[hidelinks,breaklinks]{hyperref}
\usepackage{xurl}
\setmainfont{Times New Roman}
\setmonofont{Consolas}
\setCJKmainfont{SimSun}[AutoFakeBold=true,AutoFakeSlant=true]
\setCJKsansfont{SimHei}
\setCJKmonofont{FangSong}
\renewcommand{\abstract}{\begin{center}\bfseries 摘\quad 要\end{center}}
\renewcommand{\endabstract}{\par}
\makeatletter
\renewcommand{\fnum@figure}{\textbf{图\ \thefigure}}
\renewcommand{\fnum@table}{\textbf{表\ \thetable}}
\renewenvironment{thebibliography}[1]
  {\section{参考文献}\linespread{1.1}\small\list{[\arabic{enumi}]}{\settowidth\labelwidth{[#1]}\leftmargin\labelwidth\advance\leftmargin\labelsep\usecounter{enumi}\setlength\itemsep{2pt}}\sloppy\interlinepenalty10000\clubpenalty10000\widowpenalty10000\sfcode`\.=1000\relax}
  {\endlist}
\makeatother
\setlength{\parindent}{2em}
\setlength{\abovedisplayskip}{5pt plus 2pt minus 1pt}
\setlength{\belowdisplayskip}{5pt plus 2pt minus 1pt}
\setlength{\textfloatsep}{12pt plus 2pt minus 2pt}
\setlength{\dbltextfloatsep}{12pt plus 2pt minus 2pt}
\renewcommand{\arraystretch}{1.2}
\raggedbottom
\title{融合日周周期与输入相关门控的轻量级电力负荷预测}
\IfFileExists{author.tex}{\input{author}}{\name{作者姓名}
\address{学校与学院\\班级与学号}
}
\hypersetup{pdftitle={融合日周周期与输入相关门控的轻量级电力负荷预测},pdfauthor={李军宏},pdfsubject={深度学习课程论文}}

\begin{document}
\maketitle
\begin{abstract}
电力负荷序列同时包含趋势变化与多尺度周期模式，预测模型需要在精度与计算开销之间进行权衡。本文以DLinear为基础，加入历史日周期、周周期预测分支，并利用输入窗口统计量生成融合权重，构建轻量级模型PeriodGatedDLinear。在Electricity小时级数据上，以336小时为历史窗口，开展24小时和96小时预测，设置DLinear、LSTM和GRU三个基线及三项模块消融，采用三个随机种子，共纳入42次独立训练结果。完整模型的MAE分别为0.1998和0.2292，相比DLinear降低6.32\%和5.10\%。消融结果中，周周期分支的平均收益更明显，输入相关门控优于静态权重，日周期分支带来较小改善。完整模型在96小时预测中包含77\,635个参数，实测GPU张量峰值显存为84.41\,MiB。结果说明，在所选基线和实验配置下，显式周期先验与轻量门控能够改善线性负荷预测。项目代码与复现实验脚本在GitHub开放获取。
\end{abstract}
\vspace{4pt}
\noindent\textbf{关键词：}电力负荷预测；时间序列；DLinear；多周期建模；动态门控

\section{引言}
电力负荷预测的任务是根据历史用电记录估计未来需求，可为用电计划与资源安排提供依据。多用户负荷具有幅值差异，序列内部还可能叠加日常活动、周内重复模式和缓慢变化的趋势。LSTNet将卷积、循环结构与自回归项结合，用于建模长短期时间模式\cite{lai2018lstnet}；DeepAR通过在相关序列之间共享循环网络参数进行概率预测\cite{salinas2017deepar}。这些研究说明，如何表达重复模式、共享统计信息以及控制模型规模，是负荷预测中的重要问题。

深度学习方法提供了多种时间依赖建模方式。LSTM通过记忆单元和门结构处理长序列\cite{hochreiter1997lstm}，GRU以较紧凑的门控状态更新实现序列编码\cite{cho2014gru}；Transformer则通过注意力机制构建位置之间的交互\cite{vaswani2017attention}。然而，复杂结构并不必然在所有预测场景下优于简单模型。DLinear基于分解与线性映射构建有竞争力的预测基线\cite{zeng2023dlinear}，提示研究者应在相同数据处理和评估条件下验证结构设计的实际价值。

对于显存约4\,GiB的单卡实验环境，仅追求较大模型会增加调试与重复实验成本。本文关注一个较具体的问题：在保留DLinear低开销特征的条件下，显式加入24步、168步周期预测，并根据当前输入决定各分支权重，能否改善电力负荷预测？固定融合无法区分当前窗口的周期一致性与波动状态，而单一周期也可能遗漏其他时间尺度的信息，因此有必要分别考察周期预测分支与输入相关门控的作用。

本文的工作包括三个方面。首先，在开源Time-Series-Library上实现由DLinear、历史日周期与周周期三个预测分支组成的模型，明确各分支的输入边界与相位对齐方式。其次，在统一的时间划分和评价流程下，与DLinear、标准LSTM和标准GRU三个基线比较。最后，设计去除日分支、去除周分支以及用静态权重替换动态门控的三项消融，同时报告参数量、训练时间和张量峰值显存。研究重点在于分析分解、周期先验与专家融合这些已有思想的组合效果及其适用条件。

\section{相关工作}
\subsection{循环建模与直接多步预测}
LSTM和GRU适于对按时间排列的观测进行递推编码，其训练成本与序列长度相关\cite{hochreiter1997lstm,cho2014gru}。LSTNet进一步利用不同结构处理局部模式与较长周期\cite{lai2018lstnet}，DeepAR则强调跨序列学习及预测分布\cite{salinas2017deepar}。本文的循环基线采用最终隐藏状态连接线性输出头，一次性输出全部未来值；其结果对应这一特定实现，不能代表所有循环预测模型的最佳性能。

直接多步预测也可建立在前馈结构上。N-BEATS采用基展开和残差连接实现时间序列预测\cite{oreshkin2020nbeats}；DLinear利用移动平均分解后分别进行线性投影\cite{zeng2023dlinear}。前者说明前馈结构具有表达复杂预测映射的潜力，后者为分析线性映射与附加模块的贡献提供了清晰起点。本文选择DLinear作为主体，以便观察额外周期建模与融合机制带来的变化。

\subsection{注意力与多尺度时间表示}
Transformer的注意力机制为时间依赖建模提供了基础\cite{vaswani2017attention}。Informer通过稀疏注意力等设计面向长序列预测\cite{zhou2021informer}；Autoformer将分解与自相关机制结合\cite{wu2021autoformer}，FEDformer在分解结构中引入频域建模\cite{zhou2022fedformer}。这些方法从不同角度利用长期依赖和重复结构，但具体资源成本仍取决于实现和配置。

iTransformer以变量为表示单元处理多变量相关性\cite{liu2024itransformer}。TimesNet通过多周期二维变化表示组织时间模式\cite{wu2023timesnet}，TimeMixer则以分解和多尺度混合处理不同采样尺度的信息\cite{wang2024timemixer}。这些工作为多变量和多尺度预测提供了不同思路。本文以DLinear和两类经典循环网络为实验基线，重点检验显式日周分支相对于线性分支的增量贡献。

\subsection{归一化与自适应融合}
不同用户的负荷尺度不一致，训练分布与未来分布也可能存在变化。RevIN通过可逆的实例归一化与反归一化缓解分布偏移\cite{kim2022revin}。本文对全部对比模型统一使用训练段逐用户标准化，以控制用户之间的量纲差异；窗口级可逆归一化与训练段标准化属于不同机制，本文不额外加入前者。

自适应融合的思想可追溯至混合专家模型，即根据输入分配不同专家的贡献\cite{jacobs1991moe}。Temporal Fusion Transformer也包含特征选择和门控组件，用于多步预测\cite{lim2021tft}。本文采用三个低成本预测分支和一个共享门控网络，属于稠密加权融合：每次前向均计算所有分支，再根据当前历史窗口确定各路输出的相对贡献。

\section{模型方法}
\subsection{任务定义与总体结构}
设历史输入为$X\in\mathbb{R}^{B\times L\times C}$，其中$B$为批大小，$L$为历史长度，$C$为用户数；目标为$Y\in\mathbb{R}^{B\times H\times C}$。本研究取$L=336$、$C=321$，并分别训练$H=24$和$H=96$的模型。对用户$c$，仅用训练段计算均值$\mu_c$和标准差$\sigma_c$，并将该用户所有时间段的观测变换为
\begin{equation}
\widetilde x_{t,c}=(x_{t,c}-\mu_c)/\sigma_c.
\label{eq:standard}
\end{equation}
后续公式中的$X$表示该标准化后的历史窗口。模型同时产生线性、日周期和周周期三路预测，通过历史统计量得到融合系数，其结构见图\ref{fig:model}。

\begin{figure*}[t]
\centering
\resizebox{.96\textwidth}{!}{\begin{tikzpicture}[
  font=\small, >=Latex,
  box/.style={draw=black!70,rounded corners=2pt,align=center,minimum height=11mm,inner sep=3mm},
  branch/.style={box,minimum width=43mm,fill=black!3},
  flow/.style={->,line width=.65pt}]
\node[box,minimum width=56mm] (input) at (0,0) {历史负荷窗口 $X$\\$B\times336\times321$};
\node[branch] (linear) at (-5.6,-2.05) {DLinear 分支\\移动平均分解与线性映射};
\node[branch] (daily) at (0,-2.05) {日周期分支 $P=24$\\14日模板均值与残差 MLP};
\node[branch] (weekly) at (5.6,-2.05) {周周期分支 $P=168$\\2周模板均值与残差 MLP};
\coordinate (split) at (0,-.92);
\draw[flow] (input) -- (linear);
\draw[flow] (input) -- (daily);
\draw[flow] (input) -- (weekly);
\node[box,minimum width=51mm,fill=blue!4] (features) at (5.6,0) {每样本每用户的8维统计量\\仅来自历史窗口};
\draw[flow] (input) -- (features);
\node[box,minimum width=51mm,fill=blue!4] (gate) at (5.6,1.65) {门控 MLP 与 Softmax\\$\alpha\in\mathbb{R}^{B\times321\times3}$};
\draw[flow] (features) -- (gate);
\node[box,minimum width=64mm,fill=blue!4] (fusion) at (0,-4.15) {三个预测分支加权求和\\同一窗口内跨预测步共享门控权重};
\draw[flow] (linear) |- (fusion.west);
\draw[flow] (daily) -- (fusion);
\draw[flow] (weekly) |- (fusion.east);
\draw[flow,dashed] (gate.east) -- ++(1.15,0) |- ([yshift=2mm]fusion.east);
\node[box,minimum width=56mm] (output) at (0,-5.9) {未来负荷预测 $\widehat Y$\\$B\times H\times321$，$H\in\{24,96\}$};
\draw[flow] (fusion) -- (output);
\end{tikzpicture}
}
\caption{PeriodGatedDLinear整体结构。实线表示历史输入与预测的计算路径，虚线表示融合权重。所有分支只使用预测起点之前的观测。}
\label{fig:model}
\end{figure*}

\subsection{DLinear预测分支}
遵循DLinear的分解与映射方式\cite{zeng2023dlinear}，首先用窗口为25的移动平均分离平滑项与剩余项：
\begin{equation}
X^{\mathrm{tr}}=\operatorname{MA}_{25}(X),\quad
X^{\mathrm{se}}=X-X^{\mathrm{tr}}.
\end{equation}
移动平均在序列两端重复边界值进行填充。对每个样本和用户的时间向量分别应用两个$L\rightarrow H$线性层，参数在用户之间共享：
\begin{equation}
\widehat{\boldsymbol y}^{(0)}_{b,c}
=W_{\mathrm{tr}}\boldsymbol x^{\mathrm{tr}}_{b,c}
+W_{\mathrm{se}}\boldsymbol x^{\mathrm{se}}_{b,c}
+\boldsymbol\beta_{\mathrm{tr}}+\boldsymbol\beta_{\mathrm{se}}.
\end{equation}
该分支独立提供基于全部历史时间点的直接预测；日周分支则增加更明确的重复模式约束。

\subsection{历史日周周期预测分支}
对周期$P\in\{24,168\}$，记$n_P=\lfloor L/P\rfloor$。只保留最近$n_PP$步，按相位对齐后平均完整周期，得到模板：
\begin{equation}
T^{(P)}_{b,c,j}=\frac{1}{n_P}\sum_{i=0}^{n_P-1}
X_{b,L-n_PP+iP+j,c},
\label{eq:template}
\end{equation}
其中$j=0,\ldots,P-1$，时间下标从0开始。$L=336$时，日模板平均14个完整日周期，周模板平均2个完整周周期。模板采用相对于预测起点的相位，因此不需要真实节假日或外部日历标签。

周期均值可能无法表达最近窗口中的形状调整。为此，使用在用户之间共享的残差MLP修正模板向量：
\begin{equation}
\begin{aligned}
\boldsymbol z^{(P)}&=\mathcal D\!\left(\phi(W^{(P)}_1\boldsymbol t^{(P)}+\boldsymbol b^{(P)}_1)\right),\\
\boldsymbol u^{(P)}&=\boldsymbol t^{(P)}+W^{(P)}_2\boldsymbol z^{(P)}+\boldsymbol b^{(P)}_2.
\end{aligned}
\label{eq:periodmlp}
\end{equation}
其中$\phi$为GELU激活函数\cite{hendrycks2016gelu}，$\mathcal D$表示训练阶段概率为0.1的Dropout\cite{srivastava2014dropout}，MLP维度为$P\rightarrow32\rightarrow P$。最后一层权重与偏置初始化为0，使分支初始输出等价于历史周期均值。预测阶段通过相位索引延拓：
\begin{equation}
\widehat Y^{(P)}_{b,h,c}=U^{(P)}_{b,c,h\bmod P},
\quad h=0,\ldots,H-1.
\end{equation}
这一步直接生成未来值，既不递归使用自身预测，也不读取目标区间的真实观测。周分支在$H=96$时使用修正周模板的前96个对齐相位；日分支则将24个相位重复至完整预测区间。

\subsection{输入相关门控与预测融合}
对每个样本、每个用户提取8项历史统计量，组成$\boldsymbol s_{b,c}\in\mathbb{R}^8$，依次为全窗口均值、总体标准差、最后观测、最近24步均值、最近24步总体标准差、首尾差，以及24步和168步滞后的平均绝对差。其中滞后差定义为
\begin{equation}
d_P=\frac{1}{L-P}\sum_{t=P}^{L-1}|X_{b,t,c}-X_{b,t-P,c}|.
\end{equation}
较小的$d_P$表示当前历史窗口在该滞后下具有较强的一致性，但不保证未来仍保持相同模式。

门控MLP的维度为$8\rightarrow32\rightarrow3$，GELU后不使用Dropout，其输出经Softmax归一化：
\begin{equation}
\boldsymbol\alpha_{b,c}=\operatorname{Softmax}
\left(W_{g,2}\phi(W_{g,1}\boldsymbol s_{b,c}+\boldsymbol b_{g,1})
+\boldsymbol b_{g,2}\right).
\end{equation}
三个系数非负且和为1，最终预测为
\begin{equation}
\widehat Y_{b,h,c}=\sum_{k\in\{0,24,168\}}
\alpha_{b,c,k}\widehat Y^{(k)}_{b,h,c}.
\label{eq:fusion}
\end{equation}
同一用户在一个输入窗口内的全部预测步共享一组权重，因此该设计是输入窗口相关门控，并非逐预测步变化的门控。门控最后一层采用标准差0.01的小随机权重，偏置设为$(2,0,0)$，使初始预测倾向DLinear分支。

各预测分支和门控参数采用端到端训练，优化标准化目标上的均方误差：
\begin{equation}
\mathcal L=\frac{1}{BHC}\sum_{b=1}^{B}\sum_{h=1}^{H}\sum_{c=1}^{C}
(\widehat Y_{b,h,c}-Y_{b,h,c})^2.
\end{equation}
模型在用户之间共享分支和门控参数，不显式混合不同用户的观测。它适合检验单用户历史结构的作用，但无法直接表示用户间的空间或相关关系。

\section{实验设置}
\subsection{数据集与时间划分}
数据源追溯至UCI的Electricity负荷数据集\cite{trindade2015electricity}。本文实际使用Time-Series-Library数据发布页提供的小时级预处理CSV\cite{tsldata}，而不是直接处理UCI的370用户、15分钟原始文件。预处理版本包含26\,304个连续小时、321个负荷变量，文件大小为95\,581\,762字节，校验未发现缺失值或非有限数值。

发布CSV的时间索引起于2016年7月1日02:00，止于2019年7月2日01:00，与原始数据集名称中的年份不一致。本文保留原文件及索引，只利用相对24步和168步周期，不据此匹配真实节假日、天气或原始采集日期；321个负荷变量全部作为预测目标。

原始时间点按约70\%、10\%、20\%划分为训练、验证和测试段，数量分别为18\,412、2\,632和5\,260。验证和测试窗口可使用各自边界之前已观测的336步历史，但预测目标必须完整位于相应分段内，标准化始终只拟合训练段。表\ref{tab:windows}给出滑动窗口数，滑动步幅为1；测试评价包含最后一个不足整批的窗口批次。

\begin{table}[t]
\centering
\caption{各分段的可用滑动窗口数}
\label{tab:windows}
\vspace{4pt}
\small\setlength{\tabcolsep}{9pt}
\begin{tabular}{crrr}\toprule
预测长度 & 训练 & 验证 & 测试\\\midrule
24 & 18\,053 & 2\,609 & 5\,237\\
96 & 17\,981 & 2\,537 & 5\,165\\\bottomrule
\end{tabular}
\end{table}

\subsection{基线与训练配置}
实验基于Time-Series-Library\cite{tslcode}，底层使用PyTorch\cite{paszke2019pytorch}。对比范围包含DLinear、标准LSTM、标准GRU和本文Core模型。DLinear采用用户共享的线性层；LSTM与GRU均为2层、隐藏维度64、层间Dropout为0.1，由顶层最终隐藏状态通过线性头直接输出$H\times321$个预测值。Core使用前述日周分支与输入相关门控，全部对比模型的批大小统一为16。

所有模型使用Adam优化器\cite{kingma2015adam}，初始学习率为0.001，采用余弦学习率调整，梯度范数裁剪阈值为1，最多训练20轮。验证MSE连续3轮未改善则早停，最终重新加载验证最优参数后进行测试，测试集不参与训练期间的模型选择。不同配置的实际早停轮数可能不同，因此最终训练更新次数并不完全一致。

每种配置分别设置2021、2022、2023三个随机种子，控制随机初始化、训练顺序和Dropout等过程，报告三次运行的均值与样本标准差。本文纳入的对比实验为4模型$\times$2预测长度$\times$3种子，共24次；三项消融增加18次，完整模型复用前24次中的结果，合计42次正式运行。未将调试阶段的限批测试结果纳入表格。

全部记录来自NVIDIA GeForce RTX 3050 Laptop GPU，显存4\,GiB，运行环境为Python 3.11.5与PyTorch 2.7.1+cu118。不同模型的参数规模和早停轮数存在差异，因此资源数据用于描述当前实现的实际开销，不等同于严格控制计算预算的算法效率排名。

\subsection{评价指标与统计口径}
设某次运行共有$N$个测试窗口，标准化空间中的误差为$e_{n,h,c}=\widehat Y_{n,h,c}-Y_{n,h,c}$，定义
\begin{align}
\operatorname{MAE}&=\frac{1}{NHC}\sum_{n,h,c}|e_{n,h,c}|,\\
\operatorname{RMSE}&=\sqrt{\frac{1}{NHC}\sum_{n,h,c}e_{n,h,c}^2}.
\end{align}
两项指标均对全部测试窗口、预测步和用户进行汇总，不具有kW等物理单位。每个种子先独立计算RMSE，再对三个RMSE取平均，不以平均MSE开方替代该统计过程。对三次结果$m_i$，样本标准差为
\begin{equation}
s_m=\sqrt{\frac{1}{3-1}\sum_{i=1}^{3}(m_i-\bar m)^2},
\quad \bar m=\frac{1}{3}\sum_{i=1}^{3}m_i.
\end{equation}
标准差反映随机种子带来的波动，不是预测区间。滑动窗口之间高度重叠，本文不将其视作独立重复实验，也不据三个种子声称统计显著性。

训练时间从数据及模型准备完成后开始计时，包含训练、验证和最优检查点保存，排除最终测试。峰值显存取训练与验证期间PyTorch的\nolinkurl{max_memory_allocated}，以MiB计（$1\,\mathrm{MiB}=2^{20}$字节），仅反映张量分配峰值，不含全部驱动、CUDA上下文和桌面显示开销。

\section{实验结果与分析}
\subsection{预测精度对比}
\begin{table*}[t]
\centering
\caption{对比实验结果（3个种子的均值$\pm$样本标准差，误差越小越好；粗体为各列最优）}
\label{tab:main}
\vspace{4pt}
\small\setlength{\tabcolsep}{9pt}
\begin{tabular}{lcccc}\toprule
模型 & \multicolumn{2}{c}{$H=24$} & \multicolumn{2}{c}{$H=96$}\\
\cmidrule(lr){2-3}\cmidrule(lr){4-5}
 & MAE & RMSE & MAE & RMSE\\\midrule
DLinear & $0.2132\pm0.0005$ & $0.3340\pm0.0003$ & $0.2415\pm0.0020$ & $0.3779\pm0.0009$\\
LSTM & $0.3602\pm0.0058$ & $0.5140\pm0.0044$ & $0.3790\pm0.0066$ & $0.5415\pm0.0049$\\
GRU & $0.3645\pm0.0015$ & $0.5117\pm0.0005$ & $0.3880\pm0.0030$ & $0.5500\pm0.0017$\\
Core & $\mathbf{0.1998\pm0.0016}$ & $\mathbf{0.3195\pm0.0014}$ & $\mathbf{0.2292\pm0.0005}$ & $\mathbf{0.3656\pm0.0004}$\\
\bottomrule\end{tabular}
\end{table*}
表\ref{tab:main}显示，Core在两个预测长度下均优于DLinear。$H=24$时，MAE从0.2132降至0.1998，按未舍入均值计算相对降低6.32\%，RMSE降低4.34\%；$H=96$时，MAE从0.2415降至0.2292，相对降低5.10\%，RMSE降低3.24\%。两类误差同时改善，说明在该数据划分与训练配置下，附加周期分支与输入相关融合为线性分支提供了有用补充。

在本文选取的三个基线范围内，Core在两个预测长度下均取得最低平均MAE和RMSE。其MAE样本标准差分别为0.0016和0.0005，描述了本组三次运行的波动。上述比较仅适用于所列基线、数据划分与训练配置，不能据此推断其优于其他未纳入本文比较的预测方法。

LSTM和GRU在当前配置下的误差明显较高。其最终隐藏状态需要压缩336步、321变量的历史，再以单个输出头恢复所有预测值，这一具体结构可能形成信息表达瓶颈。相比之下，DLinear及Core保留逐用户时间投影或周期模板，具有不同的结构先验。不过，本文没有独立验证瓶颈的因果作用，也未对循环网络进行大规模超参数搜索，因此上述解释仅作为可能机制，不能据此否定循环方法在其他配置中的能力。

从24小时扩展到96小时后，各模型平均误差均增大，但仍需注意，两种预测长度对应分别训练的模型和略有差异的可用窗口集合。这里的结果用于比较两个任务的总体难度，不等同于同一模型逐步递归时的误差累积分析。

\subsection{消融设计}
消融围绕三个模块逐项修改，其他训练配置与评价流程保持一致。去除日分支时仅保留DLinear与周分支，程序名称为Weekly；去除周分支时保留DLinear与日分支，程序名称为Daily。两种情况下动态门控输出维度变为2，但8项历史输入特征保持不变。因此，消融检验的是对应周期\emph{预测分支}的贡献，并不表示完全移除了该周期相关的信息。

去除动态门控时，保留三个预测分支，使用一组全局可学习logits经Softmax形成静态融合权重，程序名称为Static。该权重对所有样本和用户相同，初始化为$(2,0,0)$；因此Static与Core的主要区别在于融合系数是否依赖当前历史输入，而非是否允许学习权重。本文没有将该项描述为完全删除融合机制，也没有混入等权平均实验。

\subsection{模块贡献分析}
\begin{table*}[t]
\centering
\caption{三项消融实验结果（3个种子的均值$\pm$样本标准差；完整模型复用对比实验结果）}
\label{tab:ablation}
\vspace{4pt}
\small\setlength{\tabcolsep}{7pt}
\begin{tabular}{lcccc}\toprule
模型变体 & \multicolumn{2}{c}{$H=24$} & \multicolumn{2}{c}{$H=96$}\\
\cmidrule(lr){2-3}\cmidrule(lr){4-5}
 & MAE & RMSE & MAE & RMSE\\\midrule
完整模型 Core & $\mathbf{0.1998\pm0.0016}$ & $\mathbf{0.3195\pm0.0014}$ & $\mathbf{0.2292\pm0.0005}$ & $\mathbf{0.3656\pm0.0004}$\\
去除日分支（Weekly） & $0.2002\pm0.0024$ & $0.3200\pm0.0021$ & $0.2304\pm0.0020$ & $0.3663\pm0.0017$\\
去除周分支（Daily） & $0.2087\pm0.0009$ & $0.3303\pm0.0006$ & $0.2356\pm0.0027$ & $0.3723\pm0.0019$\\
去除动态门控（Static） & $0.2017\pm0.0001$ & $0.3218\pm0.0001$ & $0.2331\pm0.0006$ & $0.3687\pm0.0004$\\
\bottomrule\end{tabular}
\end{table*}
表\ref{tab:ablation}中，完整模型在四个消融指标列上均取得最低平均误差。去除日分支后，24小时与96小时MAE相对完整模型分别增加0.22\%和0.50\%；去除周分支后分别增加4.48\%和2.77\%；将动态门控替换为静态权重后分别增加0.98\%和1.69\%。这些比例均由未舍入的三种子均值计算。

周分支的平均贡献更明显。其模板可保留相同周相位下不同日期的形状，而仅有24步模板的分支需要在未来多日重复同一日形状。这一结构差异与周分支更有益的结果相符。不过，实验未提供真实工作日与周末标签的独立验证，不能将误差变化直接归因于已识别出的具体日历事件。

日分支带来的改善较小，且去除日分支后的部分种子表现可优于对应的完整模型。例如在96小时、种子2023时，Weekly的MAE为0.2288，低于Core的0.2294。这表明“完整模型平均最优”不等于每次运行均占优。结合仅三个种子和模型参数量的变化，现有证据应解释为日预测分支具有有限的平均增益，而非已证明其在所有窗口和用户上不可替代。

静态与动态融合的比较表明，固定的可学习权重未达到输入相关权重的平均预测精度，支持根据局部历史状态调整分支贡献的设计。与此同时，该替换还改变了门控部分的参数量；在96小时设置下，Static为77\,251个参数，Core为77\,635个参数。因此，这不是严格参数匹配的因果消融，结论仍需限制在当前的具体结构与训练协议内。

\subsection{资源开销与精度权衡}
\begin{table*}[t]
\centering
\caption{对比模型的实测资源开销（训练时间为3个种子的均值$\pm$样本标准差；显存为GPU张量峰值）}
\label{tab:efficiency}
\vspace{4pt}
\small\setlength{\tabcolsep}{7pt}
\begin{tabular}{lrrrrrr}\toprule
 & \multicolumn{3}{c}{$H=24$} & \multicolumn{3}{c}{$H=96$}\\
\cmidrule(lr){2-4}\cmidrule(lr){5-7}
模型 & 参数量 & 时间/s & 显存/MiB & 参数量 & 时间/s & 显存/MiB\\\midrule
DLinear & 16,176 & $163.2\pm44.6$ & 57.21 & 64,704 & $107.5\pm41.5$ & 72.14\\
LSTM & 633,112 & $167.9\pm35.0$ & 109.08 & 2,135,392 & $221.7\pm55.5$ & 139.93\\
GRU & 600,024 & $158.0\pm54.2$ & 113.30 & 2,102,304 & $167.7\pm49.7$ & 142.59\\
Core & 29,107 & $236.1\pm90.9$ & 64.52 & 77,635 & $290.4\pm50.8$ & 84.41\\
\bottomrule\end{tabular}
\end{table*}
表\ref{tab:efficiency}表明，Core的参数量从24小时任务的29\,107增至96小时任务的77\,635，主要增长来自DLinear的$L\rightarrow H$时间映射；周期模板修正与三路门控本身不随$H$增加参数。96小时任务中，LSTM与GRU的参数量均超过210万，Core的参数量小于二者的4\%；其GPU张量峰值为84.41\,MiB，低于两类循环基线的139.93\,MiB和142.59\,MiB。

当前设置下，Core在24小时与96小时任务上的平均训练时间分别为236.1秒和290.4秒，均高于表中三个基线。参数量较少并不必然带来更短的训练时间：Core需要计算多个分支，且其实际早停轮数与基线不同。因此，资源比较应同时考察参数量、实测显存和训练时间，不能仅以其中一项代替总体计算成本。

相对于DLinear，Core增加周期MLP和门控，参数、显存与实际训练时间均有上升。因此，本文方法并未在所有资源指标上优于最简单的线性基线，其价值是以额外但仍较低的实测开销换取预测误差下降。应根据使用场景对精度、训练成本与结构复杂度的需求选择模型，而非仅依据单个指标判断优劣。

\subsection{局限性与可复现性}
本研究只在一个公开预处理数据集、一个历史长度和两个预测长度上进行实验，未检验跨数据集迁移、缺失数据或异常时段鲁棒性。未加入天气与节假日等外生变量，也未评价预测不确定性和实际调度成本。小时级预处理数据的时间索引差异进一步限制了基于真实日历的解释。上述边界意味着，本文结论适用于当前基准预测设置，不能直接扩展为生产环境中的可靠性保证。

本文仅选择DLinear和两类经典循环网络进行实证对比，未覆盖其他预测模型，结论应限定在该比较范围内。实验未进行模型专属的大规模调参，参数量和训练终止轮数也并非完全一致。模块消融会同时改变模型容量，缺少参数匹配对照；三个随机种子不足以支持强统计推断。特别是日分支的收益较小，后续可通过更多种子及固定容量的替代分支检验其稳健性。

代码基于Time-Series-Library提交\nolinkurl{4e938a176710}扩展，DLinear沿用上游实现，LSTM、GRU与组合模型为本项目补充。项目代码、数据校验脚本、实验入口和论文表格生成程序开放于\url{https://github.com/lijunhong714/deeplearning_exam}。仓库保留上游许可证与来源说明，并提供固定随机种子的完整运行命令。

每次正式运行保存参数配置、逐轮损失、验证最优轮次、评价指标和源码指纹。数据文件以SHA-256校验，前16位为\nolinkurl{7e45845d54c5219b}，完整值保存在项目数据清单。测试误差通过逐批累加计算，避免将全部重叠预测张量保存在内存中。论文表格由选定的42次正式结果自动汇总，可追溯至各次运行的记录，而非根据打印后的四舍五入值再次计算。

\section{结论}
本文在DLinear基础上加入历史日周周期预测分支和输入相关门控，构建并实现了轻量级负荷预测模型。本文纳入的Electricity数据42次正式运行表明，该组合在24小时和96小时预测中均改善了DLinear的MAE与RMSE。三项消融表明，周周期分支具有更明显的平均贡献，动态融合优于静态可学习权重，日分支的额外收益相对有限。

相对于所选循环基线，本文模型具有较少的参数和较低的张量峰值显存，但其训练时间仍高于三个基线，体现出预测精度与训练开销之间的权衡。代码已通过公开仓库提供，便于核对模型结构和复现实验。未来可通过扩大对比范围、多数据集验证、参数匹配消融、用户间关系建模以及受约束的逐预测步门控，进一步明确该结构的适用条件。

\bibliographystyle{IEEEbib}
\bibliography{references}
\end{document}
